\section{Discussion}

\subsection{Why Baseline Matches Transformer}

The surprising parity between per-layer statistical aggregation (baseline) and cross-layer transformer-based token embeddings (both 80\% $\pm$ 4.2\% test accuracy) has three potential explanations:

\subsubsection{Dataset Size Effect}

With only 70 training samples, the task may not require deep feature learning. Architecture family classification is coarse-grained—distinguishing ResNet from ViT based on presence/absence of attention layers, not subtle variations within families. Per-layer statistics (mean, std, norm) capture these macro-level differences:

\begin{itemize}
\item \textbf{ResNet}: BatchNorm shift/scale parameters create bimodal weight distributions
\item \textbf{ViT}: Attention weight matrices exhibit low-rank structure (sparse patterns)
\item \textbf{EfficientNet}: Depthwise separable convolutions produce factorized kernel statistics
\item \textbf{ConvNeXt}: Layer-scaled residuals show controlled weight norms
\end{itemize}

These family signatures are detectable via shallow aggregation. Larger datasets (500+ models) may reveal within-family variations requiring cross-layer reasoning.

\subsubsection{Task Complexity}

Architecture family classification may favor local patterns over global structure. Consider harder tasks:

\begin{itemize}
\item \textbf{Property prediction} (test accuracy regression): May require modeling cross-layer interactions (residual flow, gradient propagation)
\item \textbf{Backdoor detection}: May need global anomaly detection (backdoor triggers span multiple layers)
\item \textbf{Model merging}: May require understanding layer-to-layer dependencies (mergeable vs critical layers)
\end{itemize}

Our validation establishes tokenization baseline performance. Harder tasks could reveal transformer advantages.

\subsubsection{Family-Level Separability}

Strong discriminability at layer level reduces need for global reasoning. Per-family accuracy shows:

\begin{itemize}
\item \textbf{ViT most distinct}: Attention mechanism provides clear signal (85\% transformer accuracy)
\item \textbf{ResNet-ConvNeXt overlap}: Both use residual connections (80\% baseline, 80\% transformer)
\item \textbf{EfficientNet intermediate}: Shares convolutions with ResNet but has unique scaling (80\% baseline, 75\% transformer)
\end{itemize}

Layer-wise patterns dominate family discrimination. Cross-layer attention may help only when families share similar layer-level statistics but differ in architectural composition (e.g., block ordering, connection patterns).

\subsection{When Does Cross-Layer Attention Help?}

Our results suggest cross-layer modeling provides benefits when:

\textbf{Dataset Scale Increases}: Larger model zoos (1000+ models) may exhibit within-family variations requiring fine-grained analysis. Transformer capacity (2M params) is underutilized on 70 samples.

\textbf{Task Complexity Grows}: Property prediction (continuous values), backdoor detection (rare anomalies), or model synthesis (generative modeling) may demand richer representations than family classification.

\textbf{Layer Interactions Matter}: Tasks requiring understanding of architectural composition (e.g., predicting which layers are mergeable, identifying critical layers for pruning) need cross-layer reasoning that statistical aggregation cannot provide.

\textbf{Evidence from Training Dynamics}: Transformer achieved perfect training accuracy (100\%) while baseline plateaued at 98.57\%, suggesting higher representational capacity exists but is not needed for test generalization on this dataset.

\subsection{Implications for Weight-Space Learning}

\subsubsection{Tokenization is Viable}

80\% $\pm$ 4.2\% test accuracy (vs 60\% gate, 25\% random) confirms that layer-wise processing preserves sufficient structural signal for property prediction. Future work can build on this foundation with confidence that tokenization does not lose critical information.

\subsubsection{Baselines Matter}

Simple statistical features provide strong comparison points. Researchers proposing novel weight-processing architectures should establish per-layer statistics baseline to isolate gains from architectural inductive biases vs feature quality.

\subsubsection{Capacity Calibration Needed}

Transformer overfitting (100\% train, 73.33\% val) highlights model-data mismatch. Weight-space learning datasets are small compared to vision/NLP corpora. Architecture choices must account for limited sample sizes:

\begin{itemize}
\item \textbf{Small datasets (100 models)}: Per-layer statistics or shallow transformers
\item \textbf{Medium datasets (500-1000 models)}: Regularized transformers with dropout, weight decay
\item \textbf{Large datasets (5000+ models)}: Full-capacity transformers with cross-layer attention
\end{itemize}

\subsubsection{Task-Specific Architecture Selection}

Our results suggest a decision framework:

\begin{table}[h]
\centering
\caption{Recommended approach by task type}
\begin{tabular}{lll}
\hline
\textbf{Task Type} & \textbf{Recommended Approach} & \textbf{Rationale} \\
\hline
Family classification & Per-layer statistics & Local patterns dominate \\
Property prediction & Transformer (regularized) & Cross-layer interactions likely matter \\
Backdoor detection & Transformer + anomaly detection & Global anomaly patterns \\
Model synthesis & Transformer (large) & Generative modeling requires rich representations \\
\hline
\end{tabular}
\end{table}

\subsection{Limitations}

\subsubsection{Dataset Limitations}

\textbf{Small scale}: 100 models (70 train) limits conclusions about transformer benefits. Cross-layer attention advantages may emerge on larger datasets.

\textbf{Family imbalance}: 25 models per family provides balanced splits but limited within-family diversity. Larger zoos enable finer-grained family subtype discrimination.

\textbf{Vision models only}: All models are image classifiers (ResNet, ViT, EfficientNet, ConvNeXt). Generalization to NLP transformers, RL policies, or diffusion models is unknown.

\textbf{ImageNet pretraining}: All models pretrained on same dataset. Weight patterns may differ for models trained on medical images, satellite imagery, or other domains.

\subsubsection{Methodological Limitations}

\textbf{Gate threshold calibration}: 60\% threshold set without pilot study. Higher thresholds (e.g., 70-80\%) would increase rigor but risk rejecting valid tokenization strategies.

\textbf{Single task evaluation}: Architecture family classification is coarse-grained. Validation on property prediction, backdoor detection, or model merging tasks would strengthen claims.

\textbf{No within-family analysis}: We classify ResNet vs ViT but do not distinguish ResNet-50 vs ResNet-101 or ViT-Base vs ViT-Large. Finer-grained tasks may require cross-layer reasoning.

\textbf{Baseline capacity mismatch}: 10$\times$ parameter difference (200K vs 2M) confounds comparison of tokenization strategies (statistical aggregation vs learned embeddings) with model capacity effects. A capacity-matched transformer (200K params) would better isolate tokenization quality.

\textbf{Single random seed}: All experiments use seed 42. Train/test split variance is not estimated. Confidence intervals ($\pm$4.2\%) reflect bootstrap resampling variance only, not split sensitivity.

\textbf{PoC GNN not tested}: Hypothesis h-m1 (equivariant GNN for local patterns) failed gate due to implementation shortcuts (linear layers vs true EGNN, 30 models, 10 epochs). Conclusions about complementarity are unavailable.

\subsubsection{Interpretation Cautions}

\textbf{Correlation vs causation}: Baseline-transformer parity may reflect dataset properties (family separability) rather than fundamental limits of cross-layer attention. Different datasets may show different patterns.

\textbf{Overfitting analysis}: Transformer's 100\% training accuracy suggests memorization. With more data, gap between baseline and transformer may widen.

\textbf{Tokenization strategy}: We use flatten + pad + per-layer z-score normalization. Alternative strategies (graph-based tokenization, hierarchical encodings, learned embeddings) may alter baseline-transformer comparison.

\subsection{Open Questions for Future Work}

\begin{enumerate}
\item \textbf{At what dataset size does transformer outperform baseline?} Systematic scaling study (100 $\rightarrow$ 1000 $\rightarrow$ 10000 models)

\item \textbf{Do cross-layer dependencies matter for property prediction?} Regression tasks (test accuracy, FLOPs) vs classification (family)

\item \textbf{Can GNNs capture complementary local patterns?} Proper EGNN implementation with larger datasets

\item \textbf{Do different model families require different tokenization?} ResNet (convolutional) vs ViT (attention) may need family-specific preprocessing

\item \textbf{How does tokenization quality degrade with model scale?} Billion-parameter LLMs exceed fixed-length padding; adaptive strategies needed
\end{enumerate}

\subsection{Conclusion}

Our results validate layer-wise weight tokenization (80\% $\pm$ 4.2\% test accuracy) while revealing that simple per-layer statistical aggregation matches transformer-based token embedding performance on small datasets. This suggests weight-space learning is feasible and practical, with task complexity and dataset scale determining optimal architecture choice. Future work should focus on larger model zoos, harder downstream tasks, and systematic comparison of tokenization strategies to establish when cross-layer reasoning provides measurable benefits over local statistical aggregation.
